\section{Matrix Generators}

\Sec{self_consistency} presents a cautionary example of the potential pitfalls that can be encountered when small-value approximations are used during numerical analysis of experimental data.
%
Nevertheless, as demonstrated in \App{eigenvectors_of_projection}, first-order expansions often provide useful insights into problems that are otherwise difficult to solve.
%
Expanding matrix equations to first order is greatly simplified through the use of matrix generators.

For Jones matrices, the generators of boosts and rotations appear in their axis-angle representations \eqnsp{Boost}{and}{Rotation}, where the generators of the boosts are the Pauli matrices $\pauli{k}$ and those of the rotations are $\Ci\pauli{k}$.  For small values of the angles $\beta$ and $\phi$,
%
\begin{align}
\label{eqn:small_Boost}
\boost \approx \pauli{0} + \hat{\Pv{m}} \cdot \Pv{\sigma} \beta, \\
\label{eqn:small_Rotation}
\rotat \approx \pauli{0} + \Ci \hat{\Pv{n}} \cdot \Pv{\sigma} \phi,
\end{align}
%
and terms involving products of these angles are dropped; e.g.,
\begin{equation}
\begin{split}
\label{eqn:small_Product}
\boost \rotat & \approx \left(\pauli{0} + \hat{\Pv{m}} \cdot \Pv{\sigma} \beta \right) \left( \pauli{0} + \Ci \hat{\Pv{n}} \cdot \Pv{\sigma} \phi \right) \\
& \approx \pauli{0} + \hat{\Pv{m}} \cdot \Pv{\sigma} \beta + \Ci \hat{\Pv{n}} \cdot \Pv{\sigma} \phi
\end{split}
\end{equation}

Because Mueller matrices act directly on Stokes vectors, it is useful to develop the corresponding 
infinitesimal generators so that first-order calculations can be performed on products of 
non-depolarizing Mueller matrices.

\subsection{Boost Generators}

Consider the infinitesimal Hermitian Jones matrix generated by $\pauli{i}$ with 
diattenuation parameter $\partial\beta$,
\begin{equation}
\label{eqn:small_boost_basis_Jones}
\JM{B} = \exp(\partial\beta\pauli{i}) \approx \pauli{0} + \partial\beta \, \pauli{i},
\end{equation} 
and use \Eqn{Mueller_projection} to compute the components of the equivalent Mueller matrix,
\begin{equation}
  B_\irow^\icol = \frac{1}{2} \dc{\pauli{\irow}}{\left(\JM{B} \, \pauli{\icol}\,\JM{B}^\dagger\right)}.
\end{equation}
%
Define the transformed basis matrix
\begin{equation}
\begin{split}
\pauli{\icol}' & = \JM{B} \, \pauli{\icol} \JM{B} ^\dagger \\
  & \approx \left( \pauli{0} + \partial\beta \, \pauli{i} \right) \pauli{\icol} \left( \pauli{0} + \partial\beta \, \pauli{i} \right) \\
  & \approx \pauli{\icol} + \partial\beta (\pauli{i} \pauli{\icol} + \pauli{\icol} \pauli{i}) \\
  & = \pauli{\icol} + \partial\beta \{\pauli{i}, \pauli{\icol}\}
\end{split}
\end{equation}
%
where 
\begin{equation} \label{eqn:anticommutator}
    \{\JM{A}, \JM{B}\} \equiv \JM{AB} + \JM{BA}
\end{equation} 
is the matrix anti-commutator.
%
With reference to \Eqn{Pauli_product}, the Pauli matrices satisfy 
$\{\pauli{i}, \pauli{j}\} = 2\delta_{ij}\pauli{0}$ and 
$\{\pauli{i}, \pauli{0}\} = 2\pauli{i}$; 
therefore,
%
\begin{equation}
    \pauli{\icol} = \pauli{\icol} + 2 \partial\beta \, ( \delta_{0\icol} \pauli{i} + \delta_{i\icol} \pauli{0} ).
\end{equation}
%
Note that this representation of a boost doubles the infinitesimal diattenuation parameter 
because Mueller matrices act on second-order moments of the electric field (Stokes parameters) 
rather than field amplitudes.
%
Applying trace orthogonality \eqnp{basis_orthogonal} yields
%
\begin{equation}
B_\irow^\icol = \frac{1}{2} \dc{\pauli{\irow}}{\pauli{\icol}'}
  = \delta_{\irow\icol} + 2 \partial\beta ( \delta_{0\icol}\delta_{\irow i} + \delta_{i\icol} \delta_{\irow 0} ).
\end{equation}
%
That is,
%
\begin{equation}
    \label{eqn:small_boost_basis_Mueller}
    \MM{B} = \MM{I} + 2\partial\beta \MM{K}_i,
\end{equation}
where $\MM{I}$ is the $4\times4$ identity matrix and 
%
%
\begin{equation} \label{eqn:boost_generator_Mueller_index}
    \left\{ \MM{K}_i \right\}_\irow^\icol = \delta_{0\icol}\delta_{\irow i} + \delta_{i\icol} \delta_{\irow 0}
\end{equation}
%
%
is the generator of Lorentz boosts along the $S_i$ axis; e.g.
\begin{align}
\MM{K}_1 &= \begin{pmatrix} 0 & 1 & 0 & 0 \\ 1 & 0 & 0 & 0 \\ 0 & 0 & 0 & 0 \\ 0 & 0 & 0 & 0 \end{pmatrix}\\
\MM{K}_2 &= \begin{pmatrix} 0 & 0 & 1 & 0 \\ 0 & 0 & 0 & 0 \\ 1 & 0 & 0 & 0 \\ 0 & 0 & 0 & 0 \end{pmatrix}\\
\MM{K}_3 &= \begin{pmatrix} 0 & 0 & 0 & 1 \\ 0 & 0 & 0 & 0 \\ 0 & 0 & 0 & 0 \\ 1 & 0 & 0 & 0 \end{pmatrix}
\end{align}
%
Note that \Eqn{small_boost_basis_Mueller} has the same form as the right-hand side of
\Eqn{small_boost_basis_Jones}; indeed, $\MM{B} \approx \exp(2\partial\beta \MM{K}_i)$ and, more generally
\begin{equation}
    \label{eqn:boost_Mueller_exponential}
	\MM{B}_{\hat{\Pvs{m}}}(2\partial\beta) = \exp(2\partial\beta \, {\hat{\Pv{m}}} \cdot \Pv{K}),
\end{equation}
where $\Pv{K}=(\MM{K}_1, \MM{K}_2, \MM{K}_3)$ is the ordered triple of boost generators.
%
Furthermore, with reference to \App{example_boost}, the special case of
\begin{equation}
    \MM{B}_{\hat{\Pvs{1}}}(2\partial\beta) \approx \MM{I} + 2\partial\beta \MM{K}_1 
    = \begin{pmatrix}
    1 & 2\partial\beta & 0 & 0 \\
    2\partial\beta & 1 & 0 & 0 \\
    0 & 0 & 1 & 0 \\
    0 & 0 & 0 & 1
    \end{pmatrix}
\end{equation}
is consistent with the small-value approximation of \Eqn{Stokes_boost}.
%


\subsection{Rotation Generators}

Let $\JM{R}$ be an infinitesimal unitary Jones matrix generated by $\Ci\pauli{j}$ with 
rotation angle $\partial\phi$, 
\begin{equation} \label{eqn:small_rotation_basis_Jones}
\JM{R} = \exp(\Ci\partial\phi\,\pauli{j}) \approx \pauli{0} + \Ci \partial\phi \, \pauli{j},
\end{equation} 
and use \Eqn{Mueller_projection} to compute the components of the equivalent Mueller matrix,
\begin{equation}
  R_\irow^\icol = \frac{1}{2} \dc{\pauli{\irow}}{\left(\JM{R} \, \pauli{\icol}\,\JM{R}^\dagger\right)}.
\end{equation}
%
Define the transformed basis matrix
\begin{equation}
\begin{split}
\pauli{\icol}' & = \JM{R} \, \pauli{\icol} \JM{R} ^\dagger \\
  & \approx \left( \pauli{0} + \Ci \partial\phi \, \pauli{j} \right) \pauli{\icol} \left( \pauli{0} - \Ci \partial\phi \, \pauli{j} \right) \\
  & \approx \pauli{\icol} + \Ci \partial\phi (\pauli{j} \pauli{\icol} - \pauli{\icol} \pauli{j}) \\
  & = \pauli{\icol} + \Ci \partial\phi \left[\pauli{j}, \pauli{\icol}\right]
\end{split}
\end{equation}
%
where
\begin{equation} \label{eqn:commutator}
\left[\JM{A}, \JM{B}\right] \equiv \JM{AB} - \JM{BA}
\end{equation}
is the matrix commutator.
%
For the Pauli matrices, $\left[\pauli{j}, \pauli{k}\right] = 2 \Ci \epsilon_{jkl}\pauli{l}$ and $\left[\pauli{j}, \pauli{0}\right] = 0$; therefore,
\begin{equation}
    \pauli{\icol}' = \pauli{\icol} - 2 \partial\phi \epsilon_{j \icol l}\pauli{l}.
\end{equation}
Applying trace orthogonality yields
%
\begin{equation}
\begin{split}
R_\irow^\icol 
  &= \frac{1}{2} \dc{\pauli{\irow}}{\pauli{\icol}'} \\
  &= \delta_{\irow\icol} - 2 \partial\phi \epsilon_{j \icol l} \delta_{\irow l} \\
  &= \delta_{\irow\icol} + 2 \partial\phi \epsilon_{j \irow \icol } .
\end{split}
\end{equation}
%
That is,
%
\begin{equation}
    \label{eqn:small_rotation_basis_Mueller}
    \MM{R} = \MM{I} + 2\partial\phi \MM{J}_j,
\end{equation}
where 
%
%
\begin{equation} \label{eqn:rotation_generator_Mueller_index}
    \left\{ \MM{J}_j \right\}_\irow^\icol = \epsilon_{j \irow \icol}
\end{equation}
%
%
is the generator of rotations about the $S_j$ axis; e.g.
\begin{align}
\MM{J}_1 &= \begin{pmatrix} 0 & 0 & 0 & 0 \\ 0 & 0 & 0 & 0 \\ 0 & 0 & 0 & 1 \\ 0 & 0 & -1 & 0 \end{pmatrix}\\
\MM{J}_2 &= \begin{pmatrix} 0 & 0 & 0 & 0 \\ 0 & 0 & 0 & -1 \\ 0 & 0 & 0 & 0 \\ 0 & 1 & 0 & 0 \end{pmatrix}\\
\MM{J}_3 &= \begin{pmatrix} 0 & 0 & 0 & 0 \\ 0 & 0 & 1 & 0 \\ 0 & -1 & 0 & 0 \\ 0 & 0 & 0 & 0 \end{pmatrix}
\end{align}
%
As for \Eqn{boost_Mueller_exponential}, note that $\MM{R} \approx \exp(2\partial\phi \MM{J}_j)$ and,
in general
\begin{equation}
    \label{eqn:rotation_Mueller_exponential}
	\MM{R}_{\hat{\Pvs{n}}}(2\partial\phi) = \exp(2\partial\phi \, {\hat{\Pv{n}}} \cdot \Pv{J}),
\end{equation}
where $\Pv{J}=(\MM{J}_1, \MM{J}_2, \MM{J}_3)$ is the ordered triple of rotation generators.
%
Furthermore, with reference to \App{example_rotation}, the special case of
\begin{equation}
    \MM{R}_{\hat{\Pvs{2}}}(2\chi) \approx \MM{I} + 2\chi \MM{J}_2 = \begin{pmatrix} 1 & 0 & 0 & 0 \\ 0 & 1 & 0 & -2\chi \\ 0 & 0 & 1 & 0 \\ 0 & 2\chi & 0 & 1 \end{pmatrix}
\end{equation}
is consistent with the small-value approximation of \Eqn{Stokes_rotation}.
%

\subsection{Generator Transformation Properties}

The $4 \times 4$ Mueller matrix generators inherit the commutation relations 
and underlying algebraic structure of the $2 \times 2$ Jones matrix generators;
%
that is
%
\begin{align}
\left[\MM{J}_j, \MM{K}_k\right] & = -\epsilon_{jkl} \,\MM{K}_l \label{eqn:commutator_Boost} \\
\left[\MM{J}_j, \MM{J}_k\right] & = -\epsilon_{jkl} \,\MM{J}_l \label{eqn:commutator_Rotation} \\
\left[\MM{K}_j, \MM{K}_k\right] & = \epsilon_{jkl} \,\MM{J}_l  \label{eqn:commutator_mixed} 
\end{align}

\begin{proof}
\Eqn{commutator_Boost}, in index notation
\begin{align*}
\left[\MM{J}_j, \MM{K}_k\right]_\irow^\icol 
    & \equiv \left\{ \MM{J}_j \MM{K}_k - \MM{K}_k \MM{J}_j \right\}_\irow^\icol \\
    & \hspace{-4em} = \left\{\MM{J}_j\right\}_\irow^\kappa \left\{\MM{K}_k\right\}_\kappa^\icol 
      - \left\{\MM{K}_k\right\}_\irow^\lambda \left\{\MM{J}_j\right\}_\lambda^\icol & \peqn{matrix_product} \\
    & \hspace{-4em}= \epsilon_{j\irow\kappa} \left( \delta_{0\icol}\delta_{\kappa k} + \delta_{k\icol} \delta_{\kappa 0} \right) 
        & \peqn{boost_generator_Mueller_index} \\
    & - \left( \delta_{0\lambda}\delta_{\irow k} + \delta_{k\lambda} \delta_{\irow 0} \right) \epsilon_{j\lambda\icol}
        & \peqn{rotation_generator_Mueller_index} \\
   & \hspace{-4em} = \epsilon_{j\irow k} \delta_{0\icol} - \delta_{\irow 0} \epsilon_{jk\icol} \\
   & \hspace{-4em} = -\epsilon_{jkl} \delta_{0\icol} \delta_{\irow l} - \epsilon_{jkl} \delta_{l \icol} \delta_{\irow 0} \\
   & \hspace{-4em} = -\epsilon_{jkl} \left\{ \MM{K}_l \right\}_\irow^\icol & \peqn{boost_generator_Mueller_index}
\end{align*}
\end{proof}

To prove \Eqns{commutator_Rotation}{and}{commutator_mixed}, requires
$$ \epsilon_{ijk} \epsilon_{imn} = \delta_{jm}\delta_{kn} - \delta_{jn}\delta_{km}. $$

\begin{proof}
\Eqn{commutator_Rotation}, in index notation
\begin{align*}
\left[\MM{J}_j, \MM{J}_k\right]_\irow^\icol 
    & \equiv \left\{ \MM{J}_j \MM{J}_k - \MM{J}_k \MM{J}_j \right\}_\irow^\icol \\
    & \hspace{-4em} = \left\{\MM{J}_j\right\}_\irow^\kappa \left\{\MM{J}_k\right\}_\kappa^\icol 
      - \left\{\MM{J}_k\right\}_\irow^\lambda \left\{\MM{J}_j\right\}_\lambda^\icol & \peqn{matrix_product} \\
    & \hspace{-4em} = \epsilon_{j\irow\kappa} \epsilon_{k\kappa\icol} -  \epsilon_{k\irow\lambda} \epsilon_{j\lambda\icol}
        & \peqn{rotation_generator_Mueller_index} \\
    %% & \hspace{-4em} = \epsilon_{j\irow\kappa} \epsilon_{\icol k\kappa} - \epsilon_{\irow k\lambda} \epsilon_{j \icol \lambda}
    %%     & \peqn{rotation_generator_Mueller_index} \\
    %% & \hspace{-4em} = \delta_{j\icol}\delta_{\irow k} - \delta_{jk}\delta_{\irow\icol}
    %% - \delta_{\irow j}\delta_{k\icol} + \delta_{\irow\icol}\delta_{kj} \\  
    & \hspace{-4em} = \delta_{j\icol}\delta_{\irow k} - \delta_{\irow j}\delta_{k\icol}  \\           
    & \hspace{-4em} = - \epsilon_{jkl}  \epsilon_{l \irow\icol}   \\           
   & \hspace{-4em} = - \epsilon_{jkl} \left\{ \MM{J}_l \right\}_\irow^\icol & \peqn{rotation_generator_Mueller_index}
\end{align*}
\end{proof}

Furthermore, in the following proof, note that $$\delta_{0i}=\delta_{0j}=\delta_{0k}=0.$$

\begin{proof}
\Eqn{commutator_mixed}, in index notation
\begin{align*}
\left[\MM{K}_j, \MM{K}_k\right]_\irow^\icol 
    & \equiv \left\{ \MM{K}_j \MM{K}_k - \MM{K}_k \MM{K}_j \right\}_\irow^\icol \\
    & \hspace{-4em} = \left\{\MM{K}_j\right\}_\irow^\kappa \left\{\MM{K}_k\right\}_\kappa^\icol 
      - \left\{\MM{K}_k\right\}_\irow^\lambda \left\{\MM{K}_j\right\}_\lambda^\icol \\
    & \hspace{-4em} = \left( \delta_{0\kappa}\delta_{\irow j} + \delta_{j\kappa} \delta_{\irow 0} \right) \left( \delta_{0\icol}\delta_{\kappa k} + \delta_{k\icol} \delta_{\kappa 0} \right) \\
    & \hspace{-3em} - \left( \delta_{0\lambda}\delta_{\irow k} + \delta_{k\lambda} \delta_{\irow 0} \right) \left( \delta_{0\icol}\delta_{\lambda j} + \delta_{j\icol} \delta_{\lambda 0} \right) \\
%
%   & \hspace{-4em} = \left( \delta_{\irow j}\delta_{0\icol}\delta_{0 k} + \delta_{\irow j}\delta_{k\icol}
%                          + \delta_{jk}\delta_{\irow 0}\delta_{0\icol} + \delta_{j0}\delta_{\irow 0}\delta_{k\icol} \right) \\
%
%   & \hspace{-3em} - \left( \delta_{\irow k}\delta_{0\icol}\delta_{0j} + \delta_{\irow k}\delta_{j\icol}
%                          + \delta_{kj}\delta_{\irow 0}\delta_{0\icol} + \delta_{k0}\delta_{\irow 0}\delta_{j\icol} \right) \\
%
%   & \hspace{-4em} =  \delta_{\irow j}\delta_{k\icol} + \delta_{jk}\delta_{\irow 0}\delta_{0\icol} - \delta_{\irow k}\delta_{j\icol} - \delta_{kj}\delta_{\irow 0}\delta_{0\icol}  \\
%
    & \hspace{-4em} = \delta_{\irow j}\delta_{k\icol} - \delta_{j\icol}\delta_{\irow k}  \\           
    & \hspace{-4em} = \epsilon_{jkl}  \epsilon_{l \irow\icol}   \\           
   & \hspace{-4em} = \epsilon_{jkl} \left\{ \MM{J}_l \right\}_\irow^\icol
\end{align*}
\end{proof}

Any ordered triple of generators that satisfy a commutation relation of the form that appears 
in \Eqns{commutator_Boost}{and}{commutator_Rotation} transforms under similarity transformations 
exactly as the components of a three-dimensional vector under ordinary rotations.  
%
Likewise, a similarity transformation of any exponential 
operator, $M = \exp(\theta \, \hat{\Pv{n}} \cdot \Pv{V})$, results in a classical geometric rotation 
of its axis, $\hat{\Pv{n}}$.

To prove this, first consider an infinitesimal rotation of a three-dimensional vector $\Pv{v}$ 
by an angle $\partial\theta$ about an axis $\hat{\Pv{e}}_j$,
\begin{equation}
    \Pv{v}' \approx \Pv{v} + \partial\theta \, \hat{\Pv{e}}_j \times \Pv{v},
\end{equation}
%
where the vector cross-product,
\begin{equation}
(\Pv{a} \times \Pv{b})_i =  \epsilon_{ijk} a_j b_k,
\end{equation}
and $(\hat{\Pv{e}}_j)_m = \delta_{jm}$.  Therefore,
\begin{equation}
    (\hat{\Pv{e}}_j  \times \Pv{v})_k
    = \epsilon_{kml} \delta_{jm} v_l
    = \epsilon_{kjl} v_l 
    = -\epsilon_{jkl} v_l,
\end{equation}
and the rotation in component form is
\begin{equation} \label{eqn:class_vec}
    v_k' \approx v_k - \partial\theta \, \epsilon_{jkl} v_l.
\end{equation}
%
With reference to \Eqn{small_rotation_basis_Mueller},
this transformation is equivalent to multiplying $\Pv{v}$ by the $3 \times 3$ rotation matrix 
$\partial\PM{R} \approx \PM{I} - \partial\theta \PM{J}_j$, where $(\PM{J}_j)_k^l = \epsilon_{jkl}$. 

Next, consider an operator $V_k$ that satisfies the commutation relation, $[J_j, V_k] = - \epsilon_{jkl} V_l$.
(Note that the dimension of the operator is irrelevant; therefore, the boldface and underlining conventions
adopted for matrices are not used.)
After a similarity transformation by the infinitesimal rotation operator, $\partial R \approx I + \partial\theta J_j$,
\begin{equation}
\begin{split}
    V_k' &= \partial R \, V_k \ \, \partial R^{-1} \\
         &\approx (I + \partial\theta J_j) V_k (I - \partial\theta J_j) \\
         &\approx V_k + \partial\theta (J_j V_k - V_k J_j) \\
         &= V_k + \partial\theta \, [J_j, V_k] \\
         &= V_k - \partial\theta \, \epsilon_{jkl} V_l.
\end{split}
\end{equation}
%
Comparing with \Eqn{class_vec}, the operator has been rotated around axis $j$ by $\partial\theta$. 

To extend this result from an infinitesimal rotation to a finite rotation angle $\theta$, 
let $R(\theta)=\exp(\theta J_j)$ denote a rotation about axis $j$ and note that 
$dR(\theta) = J_j R(\theta) d\theta$. 
Differentiating $V_k(\theta)= R(\theta) V_k R^{-1}(\theta)$ with respect to $\theta$ yields
\begin{align} \label{eqn:ode_op}
\frac{d}{d\theta} V_k(\theta)
    &= \left( \frac{d}{d\theta} R(\theta) \right) V_k R^{-1}(\theta) + R(\theta) V_k \left( \frac{d}{d\theta} R^{-1}(\theta) \right) \nonumber \\
    &= J_j R(\theta) V_k R^{-1}(\theta) - R(\theta) V_k R^{-1}(\theta) J_j \nonumber \\
    &= [J_j, V_k(\theta)] \nonumber \\
    &= - \epsilon_{jkl} V_l(\theta).
\end{align}
This linear system of coupled differential equations can be compared with that of a three-dimensional
vector $\Pv{v}(\theta) = \PM{R}(\theta)\Pv{v}$ rotating about the same axis $j$. 
Its rate of change, $d\Pv{v}(\theta) = \hat{e}_j \times \Pv{v}(\theta) d\theta$, in component form is
\begin{equation} \label{eqn:ode_vec}
    \frac{d}{d\theta} v_k(\theta) = -\epsilon_{jkl} v_l(\theta).
\end{equation}
Comparing \Eqns{ode_op}{and}{ode_vec}, the operator $V_k(\theta)$ and the vector 
components $v_k(\theta)$ satisfy identical systems of first-order linear differential equations
with identical initial conditions at $\theta = 0$ (namely $V_k(0) = V_k$ and $v_k(0) = v_k$). 
%
By the existence and uniqueness theorem for linear systems of ODEs, $V_k(\theta)$ must evolve identically to $v_k(\theta)$ for all finite angles $\theta$. Therefore, the finite similarity transformation evaluates to
\begin{equation} \label{eqn:generator_rotation}
    R(\theta) V_k R^{-1}(\theta) = \PM{R}_k^l(\theta) V_l
\end{equation}

Finally, consider a finite operator $M$ constructed via the exponential map in axis-angle representation,
\begin{equation}
    M = \exp(\theta \, \hat{\Pv{n}} \cdot \Pv{V}) = \exp\left( \theta \, n_k V_k \right).
\end{equation}
%
To compute the similarity transformation, 
$$M' = R M R^{-1},$$
first note that $R \exp(A) R^{-1} = \exp(R A R^{-1})$.

\begin{proof}
Prove that $(\JM{R} \JM{A} \JM{R}^{-1})^n = \JM{R} \JM{A}^n \JM{R}^{-1}$ via induction,
starting with the base case ($n=1$)
\begin{align*}
(\JM{R} \JM{A} \JM{R}^{-1})^1 = \JM{R} \JM{A} \JM{R}^{-1} = \JM{R} \JM{A}^1 \JM{R}^{-1} 
\end{align*}
If $(\JM{R} \JM{A} \JM{R}^{-1})^n = \JM{R} \JM{A}^n \JM{R}^{-1}$, then
\begin{align*}
(\JM{R} \JM{A} \JM{R}^{-1})^{n+1} &= (\JM{R} \JM{A} \JM{R}^{-1})^n (\JM{R} \JM{A} \JM{R}^{-1}) \\
&= \JM{R} \JM{A}^n \JM{R}^{-1} \JM{R} \JM{A} \JM{R}^{-1} \\
&= \JM{R} \JM{A}^{n+1} \JM{R}^{-1} 
\end{align*}
%
Now substitute into \eqn{matrix_exponential} to yield
\begin{align}
\exp\left(\JM{R} \JM{A} \JM{R}^{-1}\right)
& = \sum_{n=0}^\infty \frac{1}{n!} \left( \JM{R} \JM{A} \JM{R}^{-1} \right)^n \\
& = \sum_{n=0}^\infty \frac{1}{n!} \JM{R} \JM{A}^n \JM{R}^{-1} \\
& = \JM{R} \left( \sum_{n=0}^\infty \frac{1}{n!} \JM{A}^n \right) \JM{R}^{-1} \\
& = \JM{R} \exp\left(  \JM{A} \right) \JM{R}^{-1}.
\end{align}
\end{proof}
Combine with \eqn{generator_rotation} to arrive at
\begin{equation}
\begin{split}
    M' &= \exp\left( \theta \, n_k (R V_k R^{-1}) \right) \\
       &= \exp\left( \theta \, n_k \PM{R}_k^l V_l \right) 
\end{split}
\end{equation}
With reference to \Eqn{matrix_times_vector}, 
\begin{equation}
n_k \PM{R}_k^l = (\PM{R}^T \hat{\Pv{n}})_l.
\end{equation}
That is
\begin{equation}
    M' = \exp\left( \theta\, n'_l V_l \right) = \exp( \theta \, \hat{\Pv{n}}' \cdot \Pv{V} )
\end{equation}
where $\hat{\Pv{n}}' = \PM{R}^T \hat{\Pv{n}} = \PM{R}^{-1} \hat{\Pv{n}}$ is
the unit vector obtained by rotating $\hat{\Pv{n}}$ by the inverse of the rotation associated with the similarity transformation.


In summary, by applying similarity transformations and utilizing the matrix commutator 
$[J_i, V_j] = \epsilon_{ijk} V_k$, it is shown that the algebraic properties of the operator 
space perfectly map onto classical three-dimensional vector geometry. 
%
When an exponential operator $M$ is \emph{actively} rotated using a similarity transformation, 
its scalar magnitude $\theta$ (e.g., diattenuation or retardance) remains invariant, 
and the transformation is equivalent to \emph{passively} rotating its axis.
